%=============================================================================
%  chapter3.tex -- Two Clocks
%  Sections 3.1-3.3.  Uses macros from thesis_main.tex Part 2.
%=============================================================================

\chapter{Two Clocks}
\label{ch:theory}

Chapter~\ref{ch:atoms} established what can happen to a single hydrogen atom in
a plasma and how fast each process runs. This chapter asks the next question: if
we know what happens to one atom, how do we follow $10^{12}$ or more of them per
cubic centimetre, distributed over 43 energy levels, and what makes the answer
hard?

The answer to the second half is that the system has two clocks running at
once, separated by four orders of magnitude. Everything in this thesis follows
from that separation, including a result that the separation itself does not
predict, as Chapter~\ref{ch:results} shows.

\medskip
\noindent\textbf{What this chapter takes as given.} Three things are established
elsewhere and used here without re-derivation. From Chapter~\ref{ch:intro}: that
divertor conditions are inferred spectroscopically, that the Balmer
$H_\alpha/H_\beta$ ratio is the working diagnostic, and the operating range
$\Te = 1$--\SI{10}{\electronvolt}, $\nel = 10^{12}$--\SI{e15}{\per\cubic\centi\metre}
over which everything is computed. From Chapter~\ref{ch:atoms}: the 43-level
state space and the reason $\ell$ is resolved for $n \le 8$ and bundled above it;
the origin and accuracy of every rate coefficient; and the Maxwellian averaging
that turns a cross section into a rate. This chapter assumes those rates exist
and are correct, and asks only what the resulting equations do.

%-----------------------------------------------------------------------------
\section{The rate equation}
\label{sec:rate_equation}

\subsection{A balance on every level}

Consider one energy level of the hydrogen atom, say $2p$, and count atoms.
Atoms arrive in $2p$ when an electron collision excites a ground-state atom into
it, when an atom in a higher level decays into it radiatively, when a collision
knocks an atom down from above, and when a proton captures an electron directly
into it. Atoms leave $2p$ when they are excited upward, knocked downward,
ionised, or radiate to $1s$.

The rate of change of the $2p$ population is the sum of everything arriving
minus the sum of everything leaving. That is all a rate equation is.

\begin{figure}[tbp]
  \centering
  \includegraphics[width=0.86\textwidth]{fig3_4_partition}
  \caption[The collisional--radiative network and its slow/fast partition]{%
  The collisional--radiative network of atomic hydrogen and the partition used
  in Section~\ref{sec:qss}. Levels with $n \leq 8$ are resolved in orbital
  angular momentum; $n = 9$--$15$ are bundled with a statistical
  $\ell$-distribution. Proton-impact $\ell$-mixing (orange) couples sublevels
  within every resolved shell and is the dominant loss channel for the
  metastable $2s$ level, which has no electric-dipole route to the ground
  state; arrows are drawn for $n = 2$ and $n = 3$ only, the $n = 4$ sublevels
  being too closely spaced on this schematic for a legible arrow. The broad
  arrows are not elementary transitions: $\mathbf{a} =
  -\Lmat_{FF}^{-1}\Lmat_{FS}$ and $\mathbf{c} = -\Lmat_{FF}^{-1}\Svec_F$ are
  \emph{network responses}, each summing every cascade route from its
  reservoir into a given level (Section~\ref{sec:R_depends_on_b}). Vertical
  spacing is schematic and not drawn to energy scale.}
  \label{fig:network}
\end{figure}

A chemical engineer has written this equation many times. It is a species
balance on a well-mixed vessel, with one species per energy level and one
reaction per atomic process. Figure~\ref{fig:network} shows the network and the
partition that Section~\ref{sec:qss} will exploit. The only unusual feature is
the number of species, 43, and the fact that some of the ``reactions'' are
radiative and therefore proceed without a collision partner.

Write it out for a general level $p$:
\begin{equation}
  \frac{\mathrm{d}n_p}{\mathrm{d}t}
  = \underbrace{\sum_{q \neq p} C_{qp}\,\nel\,n_q}_{\text{collisions in}}
  - \underbrace{\sum_{q \neq p} C_{pq}\,\nel\,n_p}_{\text{collisions out}}
  + \underbrace{\sum_{q > p} A_{qp}\,n_q}_{\text{radiation in}}
  - \underbrace{\gamma_p\,n_p}_{\text{radiation out}}
  - \underbrace{Q_p\,\nel\,n_p}_{\text{ionisation}}
  + \underbrace{S_p\,\nion}_{\text{recombination}}
  \label{eq:level_balance}
\end{equation}

\noindent where $n_p$ is the population of level $p$ in \si{\per\cubic\centi\metre};
$C_{qp}$ is the rate coefficient for the collisional transition $q \to p$ in
\si{\cubic\centi\metre\per\second}; $A_{qp}$ is the Einstein coefficient for
spontaneous emission $q \to p$ in \si{\per\second}; $\gamma_p = \sum_{q<p}A_{pq}$
is the total radiative loss rate from $p$; $Q_p$ is the ionisation rate
coefficient out of $p$ in \si{\cubic\centi\metre\per\second}; and $S_p$ is the
recombination source per unit ion density in \si{\per\second}, given in
Eq.~\eqref{eq:source} below.

Collisional terms carry a density factor because a collision requires a
partner; radiative terms do not, because an excited atom decays whether or not
anything else is nearby. Which partner matters: electron-impact terms scale with
$\nel$, proton-impact terms with $\nion$. For the pure hydrogen plasma
considered here quasi-neutrality gives $\nion = \nel$, so both scale together and
the distinction is one of physics rather than of parametric dependence. This is the first hint of the
two clocks: radiative rates are fixed by atomic structure, collisional rates
scale with density, and the competition between them is what makes the answer
depend on where in the plasma you look.

\subsection{Matrix form}

Equation~\eqref{eq:level_balance} holds for each of the 43 levels. Stacking
them and collecting the terms that are linear in the tracked populations gives
\begin{equation}
  \boxed{\ \frac{\mathrm{d}\nvec}{\mathrm{d}t} = \Lmat\,\nvec + \Svec\,\nion\ }
  \label{eq:master}
\end{equation}

\noindent with $\nvec$ the 43-component population vector
(\si{\per\cubic\centi\metre}), $\Lmat$ the collisional--radiative rate matrix
(\si{\per\second}), and $\Svec\nion$ the recombination feed
(\si{\per\cubic\centi\metre\per\second}). The matrix elements are
\begin{align}
  L_{pq} &= C^{(e)}_{qp}\,\nel + C^{(i)}_{qp}\,\nion + A_{qp}
            \qquad (p \neq q),
  \label{eq:L_offdiag}\\
  L_{pp} &= -\left(\sum_{q\neq p}
            \left[C^{(e)}_{pq}\,\nel + C^{(i)}_{pq}\,\nion\right]
            + \gamma_p + Q_p\,\nel\right).
  \label{eq:L_diag}
\end{align}

The superscripts distinguish the two collision partners. $C^{(e)}$ denotes
electron-impact transitions and is the dominant term for changes of principal
quantum number. $C^{(i)}$ denotes proton-impact transitions, which in this model
are non-zero only for $\ell$-changing collisions within a shell
($\Delta n = 0$, $\Delta\ell = \pm1$, $n \le 8$). That term is easy to overlook
and it is not optional. Chapter~\ref{ch:atoms} shows it carries
\SI{99.92}{\percent} of the total loss rate from the metastable $2s$ level.
The reason is that $2s$ has no other fast exit: the transition $2s \to 1s$ is
forbidden to electric-dipole radiation, and no two-photon rate is included in
the present radiative dataset, so the only route out other than proton
$\ell$-mixing is electron-impact de-excitation, some three orders of magnitude
slower. Removing proton $\ell$-mixing therefore does not merely perturb the
$n=2$ manifold; it changes which process controls it. Quasi-neutrality gives
$\nion = \nel$ for the pure hydrogen plasma considered here, so both collisional
terms scale with the same density.

\noindent Note the index order: $L_{pq}$ with $p \neq q$ is the rate of the
transition $q \to p$, a \emph{gain} into $p$. Off-diagonal elements are
therefore non-negative and diagonal elements are negative. This convention is
the opposite of the one some texts use, and getting it backwards transposes the
physics; it is stated here once and used throughout.

\subsection{Why recombination is a feed stream, not a matrix element}

There is one term in Eq.~\eqref{eq:level_balance} that refuses to become a
matrix element, and understanding why is worth a paragraph.

$\Lmat$ multiplies $\nvec$. Every entry of $\Lmat\nvec$ must therefore be
proportional to the population of some \emph{tracked} level. Excitation,
de-excitation, radiative decay and ionisation all satisfy this: they take an
atom that is in the manifold and move it, or remove it. Recombination does
not. It is proportional to $\nel\nion$, the density of a species outside the
neutral manifold entirely, and so it enters as a constant feed, independent of
$\nvec$.

The chemical-engineering analogy is exact and worth holding onto for the rest of
the thesis. $\Lmat$ is the reaction-rate matrix of a network of 43 species in a
well-mixed vessel. $\Svec\nion$ is the inlet stream. And the ion population is
the reservoir feeding that inlet, a quantity this model does not evolve, which
becomes important in Section~\ref{sec:two_references} and again in
Chapter~\ref{ch:limits}.

\subsection{Units, and a source that is quadratic in density}
\label{sec:source_quadratic}

The source deserves care because its units are easy to state wrongly. Writing
it out,
\begin{equation}
  \Svec = \nel\left(\bm{\alpha}_{\mathrm{RR}}
        + \nel\,\bm{\alpha}_{\mathrm{3BR}}\right)
  \label{eq:source}
\end{equation}

\noindent where $\bm{\alpha}_{\mathrm{RR}}$ is the radiative recombination
coefficient (\si{\cubic\centi\metre\per\second}) and
$\bm{\alpha}_{\mathrm{3BR}}$ the three-body coefficient
(\si{\centi\metre\tothe{6}\per\second}). Three-body recombination requires a
second electron to carry away the excess energy, so it carries an extra factor
of $\nel$; the source as a whole is therefore \emph{quadratic} in electron
density, and $\Svec$ has units of \si{\per\second}, not
\si{\cubic\centi\metre\per\second}.

This matters physically, not just dimensionally. At $\nel = \SI{e12}{\per\cubic\centi\metre}$
the three-body channel already supplies more than \SI{10}{\percent} of the feed
into $30$ of the $36$ resolved levels, and is negligible only for $n \le 3$; at
$\nel = \SI{e15}{\per\cubic\centi\metre}$ it supplies more than 99\% into the
high-$n$ shells. The same equation therefore describes two physically different
regimes at the two ends of the density grid. It would be convenient if that
allowed the radiative and three-body datasets to be tested separately, one at
each end, and an earlier draft of this thesis proposed exactly that.
Section~\ref{sec:promise_broken} shows the measurement does not support it, the
three-body channel is already dominant at the bottom of the density grid, so
there is no radiative-only end to test against.

Dimensional check on Eq.~\eqref{eq:master}: $[\Lmat\nvec] =
\si{\per\second}\cdot\si{\per\cubic\centi\metre} =
\si{\per\cubic\centi\metre\per\second}$, and
$[\Svec\nion] = \si{\per\second}\cdot\si{\per\cubic\centi\metre}$, matching
$[\mathrm{d}\nvec/\mathrm{d}t]$. Every term is a number of atoms per unit volume
per unit time, as a balance equation must be.

%-----------------------------------------------------------------------------
\section{Conservation, and how we know the matrix is right}
\label{sec:conservation}

A $43\times43$ matrix assembled from six separate atomic datasets offers a great
many opportunities to make a mistake that produces plausible numbers. Before
using $\Lmat$ for anything, we need a test that a wrong matrix would fail.

The test comes from counting atoms. Sum Eq.~\eqref{eq:master} over all levels
with the feed switched off. Every collisional and radiative transition moves an
atom \emph{within} the manifold: it is a loss from one level and a gain to
another, so the two contributions cancel in the sum. Only ionisation genuinely
removes an atom from the bound states. Therefore
\begin{equation}
  \boxed{\ \sum_{p} L_{pq} = -Q_q\,\nel \qquad \text{for every column } q.\ }
  \label{eq:conservation}
\end{equation}

\noindent Each column of $\Lmat$ must sum to the ionisation leak from the level
that column represents, no more, no less.

The check is worth having and it is worth being precise about what it can
detect, because the answer is narrower than it looks. Each block of the
assembly writes its diagonal as minus the column sum of the off-diagonal
entries it has just written, so the identity holds by construction whatever
those entries contain. Injecting a tenfold error into one excitation rate
coefficient, deleting every de-excitation rate, and transposing the entire
input array each leave the residual unchanged. What the check does detect is an
inconsistency between the two separately stored radiative arrays, the
transition probabilities and the total decay rates: halving one without the
other moves the residual by thirteen orders of magnitude.
Chapter~\ref{ch:validation} grades it on that basis.

\subsection{A worked example}

The rule is easiest to see in a two-level system. Take $1s$ and $2p$ only, at
$\Te = \SI{3}{\electronvolt}$ and $\nel = \SI{e14}{\per\cubic\centi\metre}$, with
representative coefficients $C_{12} = \SI{1e-11}{\cubic\centi\metre\per\second}$,
$C_{21} = \SI{1e-10}{\cubic\centi\metre\per\second}$,
$A_{21} = \SI{6.27e8}{\per\second}$,
$Q_1 = \SI{1e-13}{\cubic\centi\metre\per\second}$ and
$Q_2 = \SI{1e-10}{\cubic\centi\metre\per\second}$. Then
\begin{equation}
  \Lmat =
  \begin{pmatrix}
    -1010 & 6.2701\times10^{8}\\[2pt]
    1000  & -6.2702\times10^{8}
  \end{pmatrix}\si{\per\second},
  \qquad
  \sum_p L_{pq} = \begin{pmatrix}-10 \\ -10^{4}\end{pmatrix}\si{\per\second}
  \label{eq:toy_matrix}
\end{equation}

\noindent and $Q_1\nel = \SI{10}{\per\second}$, $Q_2\nel = \SI{e4}{\per\second}$.
The rule holds.

Two features of this small matrix survive into the full one and are worth
noticing now.

First, the second column is a difference of two numbers that agree to four
significant figures: $6.2701\times10^{8}$ against $6.2702\times10^{8}$, leaving
$10^{4}$. Almost all of the $2p$ loss rate is radiative decay to the ground
state, and the physically interesting part, the ionisation leak, is four
orders of magnitude smaller. Conservation is therefore always quoted as a
\emph{relative} residual; an absolute one would be meaningless.

Second, that same near-cancellation is the seed of the two clocks. The $2p$
level empties in nanoseconds by radiating, but the atom that arrives back in
$1s$ is still in the manifold. Only the small ionisation term actually removes
it. Section~\ref{sec:two_clocks} makes this quantitative.

\subsection{The measured residual}

For the full matrix, evaluated at all 400 grid points, the largest relative
column-sum residual is
\begin{equation}
  \max_{q,\,\Te,\,\nel}
  \left|\frac{\sum_p L_{pq} + Q_q\nel}{Q_q\nel}\right|
  = 3.61\times10^{-11},
  \label{eq:conservation_measured}
\end{equation}

\noindent which is round-off for a matrix whose non-zero entries span about
twelve decades. The assembly is arithmetically sound.

It is worth being precise about what this does and does not establish. It
establishes that every transition entered as a loss also entered as a gain
somewhere. It says nothing whatever about whether the \emph{magnitude} of any
rate is correct: $\ell$-mixing, for instance, cancels exactly in the column sum
and so is invisible to this test. Checking magnitudes requires comparison
against independent sources, which is the subject of
Chapter~\ref{ch:validation}. A test that a wrong result would also have passed
is not a test.

%-----------------------------------------------------------------------------
\section{Two clocks}
\label{sec:two_clocks}

\subsection{Why the system has natural timescales at all}

Equation~\eqref{eq:master} is linear, so with $\Lmat$ and the feed held fixed its
solution is a steady state plus a sum of decaying exponentials:
\begin{equation}
  \nvec(t) = \nvec^{\mathrm{ss}}
           + \sum_{k} c_k\,\mathrm{e}^{\lambda_k t}\,\bm{v}_k,
  \qquad
  \Lmat\bm{v}_k = \lambda_k\bm{v}_k,
  \qquad
  \nvec^{\mathrm{ss}} = -\Lmat^{-1}\Svec\nion.
  \label{eq:eigen_solution}
\end{equation}

\noindent Each eigenvalue $\lambda_k$ has a negative real part; the system is
stable, as a system of decaying populations must be. Each eigenvalue contributes
a mode decaying with time constant $\tau_k = 1/|\lambda_k|$. The set of $\tau_k$
is the set of clocks the system runs on.

It is worth pausing on what Eq.~\eqref{eq:eigen_solution} has achieved, because
the whole chapter turns on it. In the populations $n_p$ the 43 balance equations
are hopelessly coupled: nothing happens to $3d$ without also happening to $2p$,
$4f$, and everything else. Diagonalising $\Lmat$ is a change of coordinates.
Instead of asking how many atoms are in each \emph{level}, ask how much of each
\emph{eigenvector} the state contains. In those coordinates the coupling
disappears entirely. The $k$th coordinate obeys $\dot{c}_k = \lambda_k c_k$,
which is the first-order decay $c_k(t) = c_k(0)\mathrm{e}^{-t/\tau_k}$ familiar
from single-step kinetics. Forty-three coupled equations become forty-three
independent ones, each with its own clock.

The eigenvectors are therefore not abstract: each is a \emph{pattern of
populations} that decays without changing shape, and any state of the system is
a superposition of those patterns fading at their own rates. Asking which
physical process a mode corresponds to is then a matter of reading which levels
carry its weight, which is what Section~\ref{sec:two_clocks} does next. One
caveat, developed in Section~\ref{sec:transient_growth}: this decomposition is
exact, but the patterns are not orthogonal, so a small state can be the sum of
two large and nearly cancelling ones.

For hydrogen at divertor conditions there are 43 of them. The striking fact,
and the one this thesis is built on, is that they are not spread out evenly.

\subsection{The spectrum}

Figure~\ref{fig:spectrum} shows all 43 relaxation times, and
Table~\ref{tab:spectrum} lists the five slowest at the benchmark point
$\Te = \benchTe$, $\nel = \benchne$.

\begin{table}[htbp]
  \centering
  \caption{The five slowest relaxation modes at the benchmark point. The gap
  between the first and second is $9982$; the largest of the remaining 41 gaps
  is the $2.80$ in this table, at $k=2$, so the slow mode is isolated from the
  fast band by a factor of $3567$.}
  \label{tab:spectrum}
  \begin{tabular}{clll}
    \toprule
    $k$ & $\lambda_k$ (\si{\per\second}) & $\tau_k$ (\si{\second})
        & $\tau_k/\tau_{k+1}$ \\
    \midrule
    0 & $-4.3999\times10^{4}$  & $2.2728\times10^{-5}$ & $9982$ \\
    1 & $-4.3919\times10^{8}$  & $2.2769\times10^{-9}$ & $1.23$ \\
    2 & $-5.3994\times10^{8}$  & $1.8521\times10^{-9}$ & $2.80$ \\
    3 & $-1.5110\times10^{9}$  & $6.6181\times10^{-10}$& $2.71$ \\
    4 & $-4.0922\times10^{9}$  & $2.4436\times10^{-10}$& -- \\
    \bottomrule
  \end{tabular}
\end{table}

There is one gap and then a continuum.

\begin{figure}[tbp]
  \centering
  \includegraphics[width=0.95\textwidth]{fig3_1_spectrum}
  \caption[Relaxation timescales of the 43 collisional--radiative modes]{%
  Relaxation timescales at the benchmark point. (a) All 43 modes. The slowest,
  $\tauslow$, is the relaxation of the ground-state population; the next,
  $\taurel$, is the slowest-decaying mode of the excited manifold, the
  bottleneck of the fast set, and their ratio is the separation $\Msep$. (b) The ratio of consecutive
  timescales, which shows that the first gap is exceptional: the largest of the
  other 41 is $2.80$, so the first exceeds every competitor by a factor of
  $3567$ and the 42 fast modes form a dense band rather than a hierarchy.}
  \label{fig:spectrum}
\end{figure}
 The slowest mode is separated from the
next by a factor of ten thousand; after that the modes crowd together, each
within a factor of about two or three of its neighbour. It is not a hierarchy of
three or four well-separated groups, as one might expect from the $n$-shell
structure. It is one very slow mode sitting far above a dense quasi-continuum
of fast ones.

\subsection{What the two modes are}

The eigenvectors identify them.

The slow mode $\bm{v}_0$ is almost entirely ground state: its $1s$ component is
$1.0000$ and every excited component is below $10^{-3}$. Its participation ratio,
\begin{equation}
  \mathrm{PR}(\bm{v}) = \frac{\left(\sum_i |v_i|^2\right)^2}{\sum_i |v_i|^4},
  \label{eq:participation_ratio}
\end{equation}
which counts the effective number of states a mode involves, unity for a mode
confined to a single level and $43$ for one spread evenly over all of them, is
$1.00$. This mode is the relaxation of the
ground-state population against the prescribed ion reservoir. It is often
called the ionisation balance, and that is the right physical picture, but the
term should be used with care here: this model holds $\nion$ fixed
(Section~\ref{sec:rate_equation}), so what relaxes is $\ngs$ alone and not a
self-consistent neutral--ion partition. Chapter~\ref{ch:limits} assesses what
that costs.

Its eigenvalue is closely related to, but not equal to, the direct ionisation
rate of the ground state. Writing
\begin{equation}
  |\lambda_0| = \eta\,Q_{1s}\,\nel ,
  \label{eq:stepwise_enhancement}
\end{equation}

\noindent the factor $\eta$ ranges from $1.16$ to $29.5$ across the grid, and
equals $3.12$ at the benchmark point. It is nowhere below unity on the grid, and
it is the \emph{stepwise-ionisation enhancement}: ground-state atoms are not only
ionised directly but also excited into the manifold and ionised from there.
Equation~\eqref{eq:reduced} shows exactly this: $\eta$ is the ratio of the
effective ionisation rate to the direct one,
\begin{equation}
  \eta = \frac{|\Omqss|}{Q_{1s}\,\nel},
  \label{eq:eta_def}
\end{equation}

\noindent the numerator being the full Schur complement and the denominator the
direct ground-state ionisation rate alone. The denominator is emphatically
\emph{not} $|\Lmat_{SS}|$: by Eq.~\eqref{eq:level_balance} the ground state has
no radiative loss and no $\ell$-mixing partner, so
$|\Lmat_{SS}| = (\sum_q C_{1s \to q} + Q_{1s})\nel$, which exceeds $Q_{1s}\nel$
by the whole of the excitation flux. Since the Schur complement returns part of
that flux to the ground state, $|\Omqss| < |\Lmat_{SS}|$ always, and the ratio
$\Omqss/\Lmat_{SS}$ is bounded above by unity, whereas $\eta$ is bounded
below by it. Note also that $\eta$ is defined through $|\lambda_0|$ while
Eq.~\eqref{eq:eta_def} is written with $|\Omqss|$; the two agree only to
$O(\Msep^{-1})$, which at the benchmark is one part in $10^{4}$ and at the
weakest grid point one part in $87$. Consistently,
$\eta$ grows with density, where collisional climbing of the ladder is easier,
and falls with temperature, where direct ionisation is already fast: at
$\Te = \SI{10}{\electronvolt}$, $\nel = \SI{e12}{\per\cubic\centi\metre}$ it is
$1.16$, and at $\Te = \SI{1}{\electronvolt}$,
$\nel = \SI{e15}{\per\cubic\centi\metre}$ it is $29.5$; almost all ionisation
there proceeds stepwise. We write
\begin{equation}
  \tauslow \equiv \frac{1}{|\lambda_0|} = \SI{22.73}{\micro\second}.
  \label{eq:tau_slow_def}
\end{equation}

\paragraph{A note on the name.} Earlier drafts of this work called
Eq.~\eqref{eq:tau_slow_def} $\tau_{\mathrm{QSS}}$, and that name was a mistake
worth explaining rather than quietly correcting. In this thesis QSS names the
elimination of the \emph{fast} excited manifold, the approximation of
Section~\ref{sec:qss}. The timescale above belongs to the \emph{slow} reservoir,
at the opposite end of the same spectrum. A reader met the sentence ``the QSS
approximation holds because $\tau_{\mathrm{QSS}}$ is long'' and had to work out
that the two occurrences of QSS referred to different subsystems. The symbol is
therefore $\tauslow$ throughout.

The obvious alternative, $\tau_{\mathrm{res}}$ for the reservoir, was rejected
for a specific reason: Sawada and Fujimoto use $\tau_{\mathrm{res}}$ for the
\emph{response} time of the excited levels \cite{SawadaFujimoto1994}, which is
this thesis's $\taurel$ and not its $\tauslow$. Chapter~\ref{ch:conclusions}
compares the two works directly, and reusing their symbol for the opposite
quantity would have made that comparison unreadable.

The next mode $\bm{v}_1$ is the excited states rearranging among themselves.
Its participation ratio is $2.64$, so it lives on two or three levels rather
than across the manifold. Its ground-state weight depends on the measure and
both are worth stating: in the raw eigenvector the $1s$ component is the largest
entry, at $1.0000$, while scaled by the equilibrium populations it is
$1.9\times10^{-4}$. The second measure is the one that matters for the
populations, and it is the reason this mode is not a ground-state mode:
\begin{equation}
  \taurel \equiv \frac{1}{|\lambda_1|} = \SI{2.277}{\nano\second}.
\end{equation}

Their ratio is the \emph{separation}
\begin{equation}
  \boxed{\ \Msep \equiv \frac{\tauslow}{\taurel} = 9982
  \quad\text{at the benchmark point.}\ }
  \label{eq:M_benchmark}
\end{equation}

The excited manifold settles ten thousand times faster than the ionisation
balance moves. Put in units a reader can hold: if the ionisation balance took a
year to equilibrate, the excited states would be done in under an hour.

\subsubsection*{Which operator $\Msep$ belongs to}

One point of notation has to be fixed here, because it is the source of an
apparent contradiction between this chapter and Chapter~\ref{ch:results}.

$\Msep$ is a property of an \emph{operator}, not of an experiment. Nothing in
$|\lambda_1|/|\lambda_0|$ refers to a disturbance, an initial condition, or a
history. Written with its argument exposed,
\begin{equation}
  \Msep[\Lmat] \equiv \frac{\tauslow[\Lmat]}{\taurel[\Lmat]}
  = \frac{|\lambda_1[\Lmat]|}{|\lambda_0[\Lmat]|},
  \qquad
  \Msep(\Te,\nel) \equiv \Msep[\Lmat(\Te,\nel)],
  \label{eq:M_functional}
\end{equation}

\noindent and Eq.~\eqref{eq:M_benchmark} is $\Msep(\benchTe,\benchne)$ while
Eq.~\eqref{eq:M_range} is the range of $\Msep(\Te,\nel)$ over the 400 grid
operators. Both are statements about matrices, and no notion of a step enters
either.

A temperature step brings a second operator into scope. Writing $\Lmat^- =
\Lmat(\Te^-,\nel)$ and $\Lmat^+ = \Lmat(\Te^-+\Delta\Te,\nel)$, the initial
populations are those of the CR equilibrium of $\Lmat^-$, but every relaxation
for $t>0$ is governed by $\Lmat^+$. The two separations
\begin{equation}
  \Msep^- \equiv \Msep[\Lmat^-], \qquad \Msep^+ \equiv \Msep[\Lmat^+]
  \label{eq:M_pre_post}
\end{equation}

\noindent are therefore both meaningful, and it is $\Msep^+$ that bounds the
post-step relaxation. At the benchmark point a step of one grid interval in
$\Te$ gives $\Msep^- = 9982$ and $\Msep^+ = 8243$. The difference is not a
disagreement between chapters: $\Msep$ falls steeply with temperature, with a
measured local slope $\mathrm{d}\ln\Msep/\mathrm{d}\ln\Te \approx -4$, so a
$4.8\,\%$ step moves it by $17\,\%$.

The alternative convention, defining $\Msep$ to mean the post-step operator
always, was considered and rejected, for a reason worth recording. The
$\Te$ grid is logarithmic with a ratio of $10^{1/49} = 1.0481$, so a step of one
grid interval from the benchmark lands exactly on the neighbouring grid point.
The post-step operator \emph{is} $\Lmat(\Te[24],\nel[5])$, one of the 400
operators whose separation Eq.~\eqref{eq:M_range} already reports. This is not
special to that step. Because $\Lmat$ is tabulated only on the grid, every step
used in this thesis is an integer number of grid intervals: the larger step
reported elsewhere as \enquote{$+\SI{0.6}{\electronvolt}$} is four intervals,
$\Te[23] \to \Te[27]$, an increment of \SI{0.6094}{\electronvolt} and a factor
$1.0481^4 = 1.2068$, for which $\Msep^+ = 4856$. \emph{Every} post-step operator
in this work is therefore also a static grid operator. Under that
convention the same matrix would carry one name in this chapter and another in
Chapter~\ref{ch:results}. Being pre-step or post-step is a role a matrix plays
in an experiment, not a property it has; the notation should not pretend
otherwise.

\subsection{The separation across the operating range}

Table~\ref{tab:spectrum} is one grid point. Across all 400,
\begin{equation}
  \SI{0.87}{\nano\second} \le \taurel \le \SI{38.9}{\nano\second},
  \qquad
  \SI{75.4}{\nano\second} \le \tauslow \le \SI{67.2}{\second},
  \qquad
  86.8 \le \Msep \le 1.73\times10^{9}.
  \label{eq:M_range}
\end{equation}

\noindent $\taurel$ varies by a factor of 45; $\tauslow$ by seven orders of
magnitude. The separation is never small: in the hottest, densest corner of the
grid, where it is weakest, the excited states are almost ninety times faster
than the ionisation balance.

The scope of that lower figure needs stating, because it is easy to quote in a
way it does not support. $\Msep = 86.8$ is the minimum over all 400 grid points,
and the point carrying it is excluded from the transient analysis of
Chapter~\ref{ch:results} by the plateau-window criterion, which requires
$\Msep > 900$ for a window to exist at all. The floor over the \emph{analysed}
set is therefore $\Msep \ge 902$, and that floor is imposed by the criterion
rather than measured. The same distinction applies to $\tauslow$, and an earlier
draft of Eq.~\eqref{eq:M_range} did not observe it: the floor
$\tauslow \ge \SI{1.18}{\micro\second}$ is the minimum over the 346 points that
satisfy $\Msep > 900$, not over the 400. Over all 400 it is
\SI{75.4}{\nano\second}, at the same $[49,7]$ corner, and that is the value
Eq.~\eqref{eq:M_range} now reports, so that all three ranges in it share one
scope. Neither number licenses a statement of the form
\emph{\enquote{the separation is at least $X$ wherever the model is used}}: the
first is a minimum over a set that includes points the analysis discards, and
the second is a restatement of the cut. Both are reported here as descriptions
of the spectrum, and the argument of Section~\ref{sec:qss} does not rest on
either.

The one-gap structure of Table~\ref{tab:spectrum} also survives across the grid,
though with far less margin than the benchmark suggests. Define the
\emph{isolation} of the slow mode as the first gap divided by the largest of the
other 41. At the benchmark it is $3567$. Its minimum over all 400 points is
$24.3$, at $\Te = \SI{10}{\electronvolt}$, $\nel =
\SI{e15}{\per\cubic\centi\metre}$, where the first gap is $86.8$ and its nearest
competitor $3.57$. A factor of 24 is still an unambiguous separation, one slow
mode and then a band, but it is $147$ times less margin than
Figure~\ref{fig:spectrum} displays, and the benchmark value must not be quoted as
though it described the grid.

That corner is worth naming, because it is where every part of this argument is
weakest at once: the minimum $\taurel$, the minimum $\tauslow$, the minimum
$\Msep$ and the minimum isolation all occur at the same grid point, $[49,7]$.
It is the single stress point of the two-clock picture. It is also, by the
plateau-window criterion, outside the set analysed in
Chapter~\ref{ch:results}, which is a limitation to state rather than a
convenience, since it means the weakest separation on the grid is never tested
against a transient. The figure script evaluates the isolation at all 400 points
and halts if it falls below ten.

$\taurel$ falls with density with a fitted logarithmic slope
$\mathrm{d}\ln\taurel/\mathrm{d}\ln\nel$ between $-0.46$ and $-0.53$ across the
temperature range, rather than the $-1$ that a purely collisional relaxation
would give. The difference is radiative decay,
which is density-independent and increasingly dominates the excited-state loss
budget as the density falls.

\subsection{Is the definition robust?}

Two definitions of $\taurel$ are available and it is worth checking they agree.
The first, used above, takes the second eigenvalue of the full $43\times43$
matrix. The second computes the slowest eigenvalue of the excited block
$\Lmat_{FF}$ alone: the relaxation of the excited manifold with the
ground state held fixed. These are different calculations and there is no
guarantee they coincide.

They agree to \SI{0.0098}{\percent} at the benchmark point and to better than
\SI{0.346}{\percent} everywhere on the grid. The excited manifold relaxes at
essentially the same rate whether or not the ground state is allowed to
participate.

It would be a mistake to read that agreement as independent evidence that the
subsystems separate, and the arithmetic says why. Block decoupling is a
perturbation in the coupling between the two blocks, and its leading correction
is of relative size $\Msep^{-1}$. At the benchmark $\Msep^{-1} =
\SI{0.0100}{\percent}$ against the measured \SI{0.0098}{\percent}. The two
agree to \SI{2}{\percent}. The disagreement between the two definitions of
$\taurel$ \emph{is} the first-order coupling correction; it is a measurement of
the spectral gap, not a second, independent demonstration of it. Quoting it as
corroboration would be the error Section~\ref{sec:conservation} warns against:
a test that a wrong result would also have passed.

Read correctly it is still worth having, and is a sharper statement than the
one it replaces. Across the grid the product $\Msep \times$ (relative
disagreement) runs from $0.074$ to $7.94$ with median $0.494$, order unity as
first-order perturbation theory requires, with a coefficient that varies by
two orders across the operating range. The two definitions differ by exactly
the amount the block structure predicts, which confirms the perturbative
picture rather than assuming it. The formal argument for separability is made
in Section~\ref{sec:qss} and rests on the Schur complement, not on this
comparison.

\subsection{What the separation does and does not license}

A separation of $10^{4}$ invites an obvious simplification: treat the excited
states as instantaneously equilibrated, and solve only for the slow subsystem.
That is the quasi-steady-state approximation, and Section~\ref{sec:qss} derives
it properly.

Two warnings, both of which matter later.

First, $\Msep$ compares two \emph{rates}. Chapter~\ref{ch:results} measures a
\emph{distance}, how far a diagnostic reading sits from the truth. Knowing how
quickly a system recovers from a displacement tells you nothing about how far it
was displaced, and these are independent quantities.

Second, eigenvalues describe the long-time behaviour of each mode in isolation.
The matrix $\Lmat$ is strongly asymmetric: radiative decay downward runs at
$\sim\SI{6.3e8}{\per\second}$ against collisional excitation upward at
$\sim\SI{e4}{\per\second}$ at the benchmark density, some five orders of
magnitude. Asymmetry of that size makes non-normality plausible. It does
not by itself prove it: a non-symmetric matrix may still be normal and have
orthogonal eigenvectors. The decisive test is the numerical one of
Section~\ref{sec:transient_growth}, where the numerical abscissa is found to be
positive at every grid point while every eigenvalue is negative. For a
non-normal matrix the modal
amplitudes can be large and opposing, and a perturbation can grow before it
decays even though every eigenvalue is negative. The spectrum bounds the
asymptotic decay; it does not bound the finite-time error.

Both warnings point the same way: the two clocks tell us the system \emph{can}
be simplified, and they say nothing about what the simplification costs. Working
that out is the business of the rest of this chapter.

%-----------------------------------------------------------------------------
\section{The quasi-steady-state approximation}
\label{sec:qss}

\subsection{The shortcut}

A separation of $10^{4}$ invites an obvious move. If the excited states settle
in nanoseconds while the ionisation balance takes microseconds, then on any
timescale relevant to the second, the first has already finished. So do not
solve for the excited states at all: assume they have reached whatever
configuration the current conditions demand, and solve only for the slow
subsystem.

This is the quasi-steady-state approximation, and every practical
collisional--radiative model uses some form of it. Without it, a plasma
simulation that resolves \si{\nano\second} timesteps in order to follow the
excited states cannot afford to run for the \si{\milli\second} an experiment
lasts.

The chemical-engineering version of this argument is the Bodenstein
steady-state approximation. For a reaction $A \to I \to P$ with a short-lived
intermediate, one sets $\mathrm{d}[I]/\mathrm{d}t = 0$ and solves algebraically
for $[I] = (\text{production})/(\text{loss})$. The intermediate stops being a
differential unknown and becomes slaved to whatever $[A]$ currently is. What
follows is the same move applied to 42 intermediates at once, coupled to each
other.

\subsection{Partition and elimination}

Split the 43 levels into a \emph{slow} subspace $\slow = \{1s\}$ and a
\emph{fast} subspace $\fast$ containing the 42 excited levels, and partition
Eq.~\eqref{eq:master} conformably:
\begin{equation}
  \frac{\mathrm{d}}{\mathrm{d}t}
  \begin{pmatrix}\nS\\[2pt]\nF\end{pmatrix}
  =
  \begin{pmatrix}
    \Lmat_{SS} & \Lmat_{SF}\\[2pt]
    \Lmat_{FS} & \Lmat_{FF}
  \end{pmatrix}
  \begin{pmatrix}\nS\\[2pt]\nF\end{pmatrix}
  +
  \begin{pmatrix}\Svec_S\\[2pt]\Svec_F\end{pmatrix}\nion .
  \label{eq:partitioned}
\end{equation}

Section~\ref{sec:two_clocks} justifies this particular partition rather than
some other: the slow eigenvector is $1.0000$ ground state with every excited
component below $10^{-3}$, so the slow mode and the ground state are the same
object to the precision that matters.

Setting $\mathrm{d}\nF/\mathrm{d}t = 0$ in the lower block gives an algebraic
system for the excited populations,
\begin{equation}
  \boxed{\ \nF^{\mathrm{QSS}}
  = -\Lmat_{FF}^{-1}\left(\Lmat_{FS}\,\nS + \Svec_F\,\nion\right).\ }
  \label{eq:qss_solution}
\end{equation}

Read the right-hand side as a sentence. The bracket is everything arriving in
the excited manifold from outside it: $\Lmat_{FS}\nS$ is excitation up from the
ground state, $\Svec_F\nion$ is recombination down from the ions. The factor
$-\Lmat_{FF}^{-1}$ divides by the loss network. It is a matrix rather than a
number only because a level's losses feed other levels, which have their own
losses, and so on; the division is over the whole network at once. For a
single Bodenstein intermediate it would collapse to $1/k_{\mathrm{loss}}$.

Two things make Eq.~\eqref{eq:qss_solution} usable. $\Lmat_{FF}$ is invertible
because it is a strictly column-diagonally-dominant matrix with negative
diagonal: every excited level loses population to the ground state or to the
continuum, so each column sum is strictly negative rather than zero, and the
Levy--Desplanques theorem applies. Physically, there are no perfectly trapped
states within the excited manifold. Its condition number lies between
$1.48\times10^{3}$ and $1.74\times10^{5}$ across the grid, so the inversion is
numerically well behaved despite the twelve decades spanned by the entries.
Both the range and the condition number are produced by
\texttt{verify\_operator\_conditioning.py} and stamped in
\texttt{validation/operator\_conditioning/}.

\subsection{The reduced dynamics}

Substituting Eq.~\eqref{eq:qss_solution} back into the slow block eliminates
$\nF$ entirely:
\begin{equation}
  \frac{\mathrm{d}\nS}{\mathrm{d}t}
  = \underbrace{\left(\Lmat_{SS} - \Lmat_{SF}\Lmat_{FF}^{-1}\Lmat_{FS}\right)}_{\Omqss}\nS
  + \left(\Svec_S - \Lmat_{SF}\Lmat_{FF}^{-1}\Svec_F\right)\nion .
  \label{eq:reduced}
\end{equation}

A 43-dimensional differential system has become a one-dimensional one. The
matrix $\Omqss$ is the \emph{Schur complement} of $\Lmat_{FF}$ in $\Lmat$, and
it has a direct physical reading: $\Lmat_{SS}$ is the direct loss from the
ground state, and $-\Lmat_{SF}\Lmat_{FF}^{-1}\Lmat_{FS}$ is the correction for
atoms that leave the ground state, wander through the excited manifold, and
either come back or are ionised on the way. It is stepwise ionisation, summed
over every possible path.

\textbf{This is the object the rest of the thesis is about.} $\Omqss$ is
precisely the effective ionisation rate that a plasma modeller using QSS would
construct, and it is the algebraic form of the effective coefficients tabulated
in frameworks such as ADAS~\cite{Summers2006}. When
Chapter~\ref{ch:results} asks what the approximation costs, it is asking what
this quantity costs.

\subsection{When is it valid?}

The elimination assumed the excited states are always at their instantaneous
steady value. The error in that assumption can be estimated directly. Write the
deviation from the instantaneous target as
$\bm{\delta}(t) = \nF(t) - \nF^{\mathrm{QSS}}(t)$ and subtract
Eq.~\eqref{eq:qss_solution} from the lower block of Eq.~\eqref{eq:partitioned}:
\begin{equation}
  \frac{\mathrm{d}\bm{\delta}}{\mathrm{d}t}
  = \Lmat_{FF}\bm{\delta}
  - \frac{\mathrm{d}\nF^{\mathrm{QSS}}}{\mathrm{d}t}.
  \label{eq:deviation}
\end{equation}

The deviation relaxes at the rate set by $\Lmat_{FF}$, whose slowest mode is
$\taurel$, and is forced by the motion of the target itself. If the driving
conditions change on a timescale $\taudrive$, then
$|\mathrm{d}\nF^{\mathrm{QSS}}/\mathrm{d}t| \sim
\|\nF^{\mathrm{QSS}}\|/\taudrive$, and balancing the two terms in
Eq.~\eqref{eq:deviation} at quasi-static equilibrium gives
$\|\bm{\delta}\| \sim \taurel\|\nF^{\mathrm{QSS}}\|/\taudrive$, that is
\begin{equation}
  \frac{\|\nF - \nF^{\mathrm{QSS}}\|}{\|\nF^{\mathrm{QSS}}\|}
  \sim \frac{\taurel}{\taudrive},
  \label{eq:qss_validity}
\end{equation}

\noindent so the approximation is accurate whenever $\taudrive \gg \taurel$. For
the self-consistent problem, where the only thing driving the excited states is
the slow subsystem itself, $\taudrive \to \tauslow$ and the condition becomes
\begin{equation}
  \Msep = \frac{\tauslow}{\taurel} \gg 1 .
  \label{eq:qss_condition}
\end{equation}

Section~\ref{sec:two_clocks} measured $\Msep = 9982$ at the benchmark and
$\Msep \ge 86.8$ at every one of the 400 grid points. \textbf{The condition is
satisfied everywhere on the grid, comfortably.} The floor carries the
qualification of Section~\ref{sec:two_clocks}: the point holding it is excluded
from the transient analysis of Chapter~\ref{ch:results}, and the floor over the
analysed set is imposed by the window criterion rather than measured.

Two qualifications must travel with that statement.

\textbf{Equation~\eqref{eq:qss_validity} is a spectral estimate, and a bound is
available.} The step from Eq.~\eqref{eq:deviation} to
Eq.~\eqref{eq:qss_validity} identified the relaxation rate of $\bm{\delta}$ with
the slowest eigenvalue of $\Lmat_{FF}$. That identification is exact only for a
normal operator. The quasi-static balance itself gives more:
$\bm{\delta} \simeq \Lmat_{FF}^{-1}\dot{\nvec}_F^{\rm QSS}$, hence
\begin{equation}
  \|\bm{\delta}\| \le \|\Lmat_{FF}^{-1}\|_2\,\|\dot{\nvec}_F^{\rm QSS}\|,
  \qquad
  \|\Lmat_{FF}^{-1}\|_2 = \frac{1}{\sigma_{\min}(\Lmat_{FF})},
  \label{eq:qss_bound}
\end{equation}

\noindent and since $\sigma_{\min} \le \min_k|\lambda_k|$ for every matrix, the
operative timescale $\|\Lmat_{FF}^{-1}\|_2$ is never shorter than the spectral
one. Measured across the grid, the ratio
$\|\Lmat_{FF}^{-1}\|_2\,|\lambda_{\min}(\Lmat_{FF})|$ lies between $1.118$ and
$1.526$, with median $1.151$: the spectral estimate understates the rigorous
timescale by at most \SI{53}{\percent}. Defining the corresponding bounded
separation $\Msep^{\rm eff} = \tauslow/\|\Lmat_{FF}^{-1}\|_2$,
\begin{equation}
  \Msep^{\rm eff} = 8659 \ \text{at the benchmark},
  \qquad
  \Msep^{\rm eff} \ge 77.3 \ \text{at all 400 grid points},
  \label{eq:M_eff}
\end{equation}

\noindent so Eq.~\eqref{eq:qss_condition} survives non-normality with a margin
reduced by \SI{13}{\percent} and not by orders of magnitude. It is worth asking
why the penalty is so small when $\kappa(\Lmat_{FF})$ reaches $1.74\times10^5$:
the answer is that the ill-conditioning lives in the fast subspace, among modes
separated by many decades in rate, and barely touches the slowest mode of
$\Lmat_{FF}$, which is the one that sets QSS validity. The operator is strongly
non-normal in a way that is almost irrelevant to this particular question.

That said, Eq.~\eqref{eq:qss_bound} is still a statement about the settled
deviation. Section~\ref{sec:two_clocks} warned that for a strongly non-normal
operator the spectrum does not control finite-time behaviour, and that warning
applies to the approach as opposed to the endpoint:
the balance in Eq.~\eqref{eq:deviation} is a \emph{quasi-static} argument, valid
once the fast transients have decayed and the deviation has settled onto the
slow manifold. It says nothing about what $\bm{\delta}$ does in the first few
$\taurel$ after a disturbance, when the non-orthogonality of the eigenvectors
can allow transient growth. Whether that matters for the quantity actually
measured is not decidable from the spectrum, and Chapter~\ref{ch:results}
settles it by integrating the full system rather than by argument.

\textbf{The condition is not sharp.} At the benchmark point
Eq.~\eqref{eq:qss_validity} gives $(\Msep^{+})^{-1} = 1.21\times10^{-4}$, the
post-step operator being the one that governs the relaxation being compared,
whereas direct integration returns a residual of order $10^{-8}$: the
approximation is four orders of magnitude better than the scaling argument
requires. The unwritten prefactor, set by how strongly the observable projects
onto the moving part of the target, is therefore small. Equation~\eqref{eq:qss_validity} predicts the
trend correctly and the magnitude not at all, and it should be quoted only for
the former.

That would appear to end the matter. It does not, and the reason is the subject
of the next section.

%-----------------------------------------------------------------------------
\section{Two reference states}
\label{sec:two_references}

% GAP-CH1: the diagnostic chain -- why a line ratio is the measured quantity
\subsection{Why ratios}

A spectrometer viewing a divertor records a line-integrated brightness. Turning
that into an absolute population requires the neutral density along the
sightline and the path length, neither of which is known to better than a factor
of a few. Practitioners therefore work with \emph{ratios}, in which both cancel.

The ratio used throughout this thesis is
\begin{equation}
  \Robs \equiv \frac{n_3}{n_4}
  = \frac{n_{3s}+n_{3p}+n_{3d}}{n_{4s}+n_{4p}+n_{4d}+n_{4f}},
  \label{eq:observable_def}
\end{equation}

\noindent the population ratio of the $n=3$ and $n=4$ shells. Up to Einstein
coefficients it is the Balmer $H_\alpha/H_\beta$ ratio, since $H_\alpha$ is the
$3\to2$ transition and $H_\beta$ the $4\to2$. The shell ratio is used here
because it is the cleaner object to derive with, not because the results are
restricted to it: Section~\ref{sec:emissivity_generalisation} shows every result
of this chapter holds for the $A$-weighted line ratio itself, and for any
non-negative linear functional of the state vector. The substitution is
therefore not free but it is also not load-bearing, and
Section~\ref{sec:balmer_equivalence} measures what it costs. The transient error
of Chapter~\ref{ch:results}, computed from the $A$-weighted line ratio and from
the shell ratio, agrees to \SI{0.014}{\percent} at the benchmark point; across
the whole grid the line-ratio value runs between $0.978$ and $0.9999$ times the
shell-ratio value, a worst case of \SI{2.2}{\percent} at the low-density edge
$\nel = \SI{e12}{\per\cubic\centi\metre}$. The residue is the $4f$ level, which
the shell sum in Eq.~\eqref{eq:observable_def} contains and $H_\beta$ cannot
see, and it is largest at the lowest density because that is where proton
$\ell$-mixing is weakest. It is convenient to normalise to the ground state as well,
writing $\rp = n_p/\ngs$, since that is the form in which the QSS solution is
usually tabulated.

\subsection{The hidden variable}

Dividing Eq.~\eqref{eq:qss_solution} through by the ground-state population
gives the QSS ratios
\begin{equation}
  \rp^{\mathrm{QSS}}
  = \frac{\left[-\Lmat_{FF}^{-1}\left(\Lmat_{FS}\nS + \Svec_F\nion\right)\right]_p}
         {\ngs}.
  \label{eq:qss_ratio}
\end{equation}

It is tempting, and in the literature common, to treat
Eq.~\eqref{eq:qss_ratio} as a function of $\Te$ and $\nel$ alone, tabulate it,
and invert a measured ratio against the table. That step deserves more scrutiny
than it usually receives.

Look at what the numerator contains. The term $\Lmat_{FS}\nS$ is proportional to
the ground-state density; the term $\Svec_F\nion$ is proportional to the ion
density. What fixes $\rp^{\mathrm{QSS}}$ is not either supply on its own but the
\emph{ratio} of the two.

It is convenient to measure that ratio against its equilibrium value. Define the
ground-state \emph{departure coefficient}
\begin{equation}
  \boxed{\ \bdep \equiv \frac{\ngs}{\Zsb{1}\,\nel\,\nion}\ }
  \label{eq:departure}
\end{equation}

\noindent where $\Zsb{1}$ is the Saha--Boltzmann coefficient of the ground
state,
\begin{equation}
  \Zsb{p}\,[\si{\cubic\centi\metre}]
  = 10^{6}\,\frac{g_p}{2g_i}
  \left(\frac{h^2}{2\pi m_e\,e\,\Te}\right)^{3/2}
  \exp\!\left(\frac{\chi_p}{\Te}\right),
  \qquad \chi_p = \frac{13.606}{p^2},
  \label{eq:saha_boltzmann}
\end{equation}

\noindent where $\Te$ and $\chi_p$ are both in \si{\electronvolt}, $e$ is the
elementary charge in \si{\coulomb} (converting \si{\electronvolt} to
\si{\joule}), the bracket evaluates in \si{\cubic\metre}, and the factor
$10^{6}$ converts to \si{\cubic\centi\metre}. Here $g_p = 2p^2$ is the
statistical weight of level $p$, $g_i = 1$ that of the bare proton, and the
factor $2$ in the denominator is the spin degeneracy of the free electron.
Temperatures are quoted in electron-volts throughout this thesis; no Boltzmann
constant appears, and $\Te$ should everywhere be read as the electron thermal
energy.

$\Zsb{p}$ is the factor for which $n_p = \Zsb{p}\nel\nion$ holds in complete
thermodynamic equilibrium, so $\bdep$ is dimensionless and equals unity there.
Departures above or below unity measure over- and under-population relative to
that thermodynamic reference; no more should be read into the sign than that.

It is worth attaching a number to $\bdep$ immediately, because the name
``departure coefficient'' suggests a small correction and it is nothing of the
kind. Evaluated on the collisional--radiative steady state, $\bdep$ runs from
$76.6$ at $[0,7]$, the coldest and densest corner of the grid, to
$2.683\times10^{5}$ at $[49,0]$, the hottest and thinnest, and equals $956$ at
the benchmark point. The ground state is over-populated relative to
Saha--Boltzmann by two to five orders of magnitude throughout, which is the
statement that a plasma at these densities is far from thermodynamic
equilibrium: radiative decay returns atoms to the ground state faster than
collisions can lift them out, and where the resonance lines escape the photons
are not replaced. A quantity that varies over three decades across the
operating range is not a small correction, and treating it as fixed is the
assumption this thesis examines.

The location of those extremes is worth pausing on, because the obvious guess is
wrong and the correction matters later. One expects the largest departure where
the gas is coldest and most neutral. It is the opposite: $\bdep$ is largest in
the \emph{hottest, thinnest} corner and smallest in the coldest, densest one.
The reason is that $\bdep$ measures distance from Saha--Boltzmann, and
Saha--Boltzmann at \SI{10}{\electronvolt} predicts a nearly fully ionised gas,
so any surviving neutral population is an enormous departure from it. At
\SI{1}{\electronvolt} Saha already predicts a mostly neutral gas, so a
predominantly neutral plasma is only a modest departure. The neutral
\emph{fraction} therefore runs the other way: solving $\Lmat\nvec + \Svec = 0$
directly and restoring $\nion = \nel$, the ratio $\ngs/\nion$ falls from $40.4$
at $[0,0]$ and $20.5$ at $[0,7]$, gases that are \SIrange{95}{98}{\percent}
neutral, to $5.3\times10^{-6}$ at $[49,7]$. \textbf{A large departure
coefficient does not mean many neutrals}, and the two quantities peak at
opposite corners of the grid.

\subsubsection*{Where the optically thin assumption holds}

That last observation forces a declaration that the preceding sections left
implicit. The Einstein coefficients in Eqs.~\eqref{eq:level_balance},
\eqref{eq:L_offdiag} and \eqref{eq:L_diag} enter bare: every emitted photon is
assumed to leave the plasma. Where the ground-state density is high that is
false: a Lyman-$\alpha$ photon is reabsorbed, the atom returns to $2p$, and
the effective $A_{2p \to 1s}$ is reduced by the population escape factor
$\Theta_P$ of Holstein~\cite{Holstein1947}.

Whether this matters is a question about $\ngs$, and the numbers above make it
answerable. Taking a homogeneous slab of thickness $D$ and a Doppler-broadened
Lyman-$\alpha$ line at $T_{\rm at} = \Te$:

\begin{center}
\begin{tabular}{lrrrrr}
  \toprule
  point & $\Te$ (eV) & $\nel$ (\si{\per\cubic\centi\metre}) &
  $\ngs$ (\si{\per\cubic\centi\metre}) & $\tau_c$ & $\Theta_P$ \\
  \midrule
  $[49,0]$ & 10 & $10^{12}$ & $1.1\times10^{7}$ & $9\times10^{-7}$ & $1.000$ \\
  $[49,7]$ & 10 & $10^{15}$ & $5.3\times10^{9}$ & $5\times10^{-4}$ & $1.000$ \\
  $[23,5]$ & 2.95 & $1.4\times10^{14}$ & $1.2\times10^{11}$ & $0.020$ & $0.986$ \\
  $[0,0]$  & 1 & $10^{12}$ & $4.0\times10^{13}$ & $11.1$ & $0.032$ \\
  $[0,7]$  & 1 & $10^{15}$ & $2.0\times10^{16}$ & $5.6\times10^{3}$ & $3\times10^{-5}$ \\
  \bottomrule
\end{tabular}
\end{center}

\noindent with $\tau_c$ and $\Theta_P$ quoted at $D = \SI{10}{\centi\metre}$;
$\Theta_P$ at $[0,0]$ is $0.48$, $0.032$ and $0.0023$ for $D = 1$, $10$ and
\SI{100}{\centi\metre} respectively, so the qualitative conclusion does not
depend on the choice, but no single figure should be quoted without it. The
model itself contains no geometry, and $D$ is an input.

The reading is clean. Over the hot half of the grid the plasma is optically
thin to Lyman-$\alpha$ to better than one part in $10^{3}$, and at the benchmark
point $\Theta_P = 0.986$: the bare-$A$ treatment is justified where the
quantitative results of this thesis are established. Along the
$\Te = \SI{1}{\electronvolt}$ edge it is not. There $\Theta_P$ falls to
$3\times10^{-5}$, the effective $A_{2p\to1s}$ collapses, and with it $\bdep^{\rm
CRE}$, the switching points $c_m/a_m$, and the ground-fed fractions of
Section~\ref{sec:ground_fed_fraction}.

\textbf{The results of this chapter are therefore stated for
$\Te \gtrsim \SI{2}{\electronvolt}$, and the cold edge of the grid is retained
for the asymptotic behaviour it exposes rather than as a prediction.} It is also
the corner where the atomic-only model is weakest for an independent reason: at
\SIrange{95}{98}{\percent} neutral and \SI{1}{\electronvolt} the gas is
molecular, and neither $\mathrm{H}_2$ nor molecular-assisted recombination
appears in $\Lmat$. Chapter~\ref{ch:limits} takes both limitations up
quantitatively.

The structural results of this chapter are statements about whatever $\Lmat$ is
supplied, and it is worth being exact about the condition they need, because it
is easy to claim too much here. The quasi-steady-state closure, the
three-argument form of $\Rqss$ and the identity
$\mathrm{d}\ln\Robs/\mathrm{d}\ln\bdep = \ffrac{3}-\ffrac{4}$ require that
$\Lmat_{FF}$, $\Lmat_{FS}$ and $\Svec_F$ contain no $\ngs$. That is what makes
Eq.~\eqref{eq:two_supplies} affine in the reservoir, and the affinity is what
makes $\ffrac{p}$ a logistic of unit width.

For a matrix with the Lyman transitions trapped at a \emph{fixed} escape factor
the condition holds, and the identities carry over with a different
$\bdep^{\rm CRE}$. For the self-consistent trapped problem it does not:
$\Theta_P$ depends on $\ngs$, so the effective transition probabilities do too,
$\bm a$ is no longer independent of the reservoir, and the unit width is lost.
The ceiling then becomes $\tanh(m|\Delta|/4)$ with $m \neq 1$ and not known in
advance. Section~\ref{sec:opacity} of Chapter~\ref{ch:limits} treats the
self-consistent case and quantifies where it matters.

A molecular reservoir breaks the condition more seriously, because
$\mathrm{H}_2$ set by recycling is a \emph{third} independent supply. Then
$n_p = a_p\ngs + c_p\nion + m_p n_{\mathrm{H}_2}$, $\ffrac{p}$ is not a
one-variable logistic, and the ceiling does not follow at all. The framework of
this chapter is two-channel by construction, and that is a statement about its
form rather than about the accuracy of its coefficients.

\textbf{$\bdep$ is a dynamical variable}, and it is worth being precise about
how it evolves. Differentiating Eq.~\eqref{eq:departure} at fixed $\nion$,
\begin{equation}
  \dot{\bdep}
  = \frac{\dot{\ngs}}{\Zsb{1}\nel\nion}
  - \bdep\left(\frac{\dot{\Zsb{1}}}{\Zsb{1}}
                + \frac{\dot{\nel}}{\nel}\right).
  \label{eq:b1_dynamics}
\end{equation}

\noindent Only the first term comes from Eq.~\eqref{eq:reduced}, which governs
$\ngs$; the second comes from the normalisation moving. This distinction matters
for the temperature-step experiment of Chapter~\ref{ch:results}. At the instant
of a step, $\ngs$ is continuous but $\Te$ is not, so $\Zsb{1}(\Te)$ jumps and
$\bdep$ changes discontinuously purely because its reference has changed. Once
$\Te$ and $\nel$ are again held fixed, $\bdep \propto \ngs$ and its subsequent
relaxation is that of the ground-state population, on $\tauslow$.

Either way, $\bdep(t)$ is not determined by the instantaneous $(\Te, \nel)$;
it carries the plasma's history. Nothing in the QSS derivation fixes it.

\subsection{Does the observable actually depend on it?}
\label{sec:R_depends_on_b}

The previous subsection shows that each \emph{individual} normalised population
depends on $\bdep$. That is not yet an argument about $\Robs$. A sceptical
reader is entitled to object: perhaps every level depends on $\bdep$ and the
dependence cancels in the ratio, leaving $H_\alpha/H_\beta$ a clean function of
$(\Te,\nel)$ after all. Ratios are chosen precisely because things cancel in
them. The objection must be answered, not assumed away.

It can be answered exactly, because Eq.~\eqref{eq:qss_solution} is linear in the
two supplies. Define
\begin{equation}
  \bm{a} \equiv -\Lmat_{FF}^{-1}\Lmat_{FS},
  \qquad
  \bm{c} \equiv -\Lmat_{FF}^{-1}\Svec_F ,
  \label{eq:ac_def}
\end{equation}

\noindent so that Eq.~\eqref{eq:qss_solution} reads
\begin{equation}
  \nF^{\mathrm{QSS}} = \bm{a}\,\ngs + \bm{c}\,\nion .
  \label{eq:two_supplies}
\end{equation}

Both vectors are dimensionless and both have a plain reading. $a_p$ is the
number of atoms that end up sitting in level $p$ per ground-state atom present:
not only those excited into $p$ directly, but those excited anywhere and
arriving at $p$ afterwards by cascade or by a further collision. The matrix
inverse is what performs that sum over every possible route. Likewise $c_p$ is
the population of $p$ per ion, from electrons captured into any level and
settling into $p$ on the way down. Neither depends on how many ground-state
atoms or ions are actually present; both are properties of $\Lmat_{FF}$,
$\Lmat_{FS}$ and $\Svec_F$ alone, and therefore functions of $(\Te,\nel)$ only.

So Eq.~\eqref{eq:two_supplies} says something simple: every excited level is
fed from two directions, and its population is the sum of two terms, one
proportional to the size of the reservoir below it and one to the reservoir
above. What distinguishes one level from another is the \emph{proportion} in
which it draws on the two.

Writing $a_m$ and $c_m$ for the sums over the $\ell$-resolved states of shell
$m$, and using $\ngs = \bdep\Zsb{1}\nel\nion$ from Eq.~\eqref{eq:departure},
\begin{equation}
  n_m^{\mathrm{QSS}}
  = \nion\left[a_m\,\bdep\,\Zsb{1}\,\nel + c_m\right]
  = a_m\,\ngs + c_m\,\nion .
  \label{eq:shell_population}
\end{equation}

\noindent The prefactor $\nion$ is common to both shells and cancels in the
ratio, as does the emitting path length. The neutral density does \emph{not}
cancel, and it is important not to read the previous sentence that way: what an
intensity ratio removes is the need to \emph{know} the absolute neutral density
and the path length in order to form it, not the dependence of the ratio on how
much neutral hydrogen is present relative to the ions. That dependence is the
subject of this chapter and it survives, inside the bracket, as $\bdep$:
\begin{equation}
  \boxed{\
  \Rqss(\bdep) = \frac{a_3\,\bdep\,\Zsb{1}\,\nel + c_3}
                      {a_4\,\bdep\,\Zsb{1}\,\nel + c_4}\ }
  \label{eq:R_of_b}
\end{equation}

Differentiating,
\begin{equation}
  \boxed{\
  \frac{\partial \Rqss}{\partial \bdep}
  = \frac{\Zsb{1}\,\nel\,\left(a_3c_4 - a_4c_3\right)}
         {\left(a_4\,\bdep\,\Zsb{1}\,\nel + c_4\right)^2}\ }
  \label{eq:dR_db}
\end{equation}

\begin{figure}[tbp]
  \centering
  \includegraphics[width=0.66\textwidth]{fig3_2_R_of_b1}
  \caption[The observable as a function of the hidden variable]{%
  The shell ratio $\Rqss$ as a function of the ground-state departure
  coefficient, at the benchmark point, from Eq.~\eqref{eq:R_of_b}. The curve is
  a one-parameter family: it asks what $\Robs$ would be if $\bdep$ took other
  values at fixed $(\Te,\nel)$. Between the two supply limits the observable
  changes by a factor of several; a table indexed on $(\Te,\nel)$ alone must
  select one point on this curve. Both axes are logarithmic, so the visible
  slope is $\mathrm{d}\ln\Robs/\mathrm{d}\ln\bdep$ directly, the quantity
  plotted in Figure~\ref{fig:ground_fed}. The marker shows the value the
  CR-equilibrium closure selects, and the shaded strip the range of
  $\bdep^{\mathrm{CRE}}$ over the whole grid.}
  \label{fig:R_of_b1}
\end{figure}

The observable is independent of $\bdep$ if and only if
\begin{equation}
  a_3c_4 = a_4c_3
  \qquad\Longleftrightarrow\qquad
  \frac{a_3}{c_3} = \frac{a_4}{c_4},
  \label{eq:cancellation_condition}
\end{equation}

\noindent that is, if and only if the $n=3$ and $n=4$ shells draw on the two
supplies in \emph{identical proportion}. There is no reason for them to. The
$n=4$ shell lies closer to the continuum and is more readily fed by
recombination; the $n=3$ shell is more readily reached by excitation from the
ground state. Chapter~\ref{ch:results} measures the difference and finds it
substantial across most of the operating range.

\subsection{The ground-fed fraction, and where the sensitivity vanishes}
\label{sec:ground_fed_fraction}

Equation~\eqref{eq:dR_db} settles whether the dependence exists. It is not,
however, the right measure of how \emph{strong} it is, because $\Robs$ itself
varies over orders of magnitude across the operating range and an absolute
derivative is not comparable between grid points. The dimensionless measure is
the logarithmic sensitivity, and it takes a form with a direct physical reading.

Define the \emph{ground-fed fraction} of shell $m$,
\begin{equation}
  \boxed{\ \ffrac{m} \equiv
  \frac{a_m\,\bdep\,\Zsb{1}\,\nel}
        {a_m\,\bdep\,\Zsb{1}\,\nel + c_m}\ }
  \label{eq:ground_fed_fraction}
\end{equation}

\noindent the share of shell $m$'s population that arrived from below, by
excitation out of the ground state, as opposed to from above by recombination.
That $\ffrac{m}$ lies in $[0,1]$ is not automatic: it requires $a_m \ge 0$ and
$c_m \ge 0$, and both are components of $-\Lmat_{FF}^{-1}$ acting on non-negative
vectors. The claim therefore needs $-\Lmat_{FF}^{-1}$ to be non-negative
\emph{elementwise}, which holds for the same reason $\Lmat_{FF}$ is invertible
(Section~\ref{sec:qss}): $-\Lmat_{FF}$ has non-positive off-diagonal entries and
strictly positive column sums, making it a non-singular M-matrix, and the
inverse of a non-singular M-matrix is elementwise non-negative. Physically,
adding atoms to one level cannot reduce the steady population of another, so the
statement is unsurprising, but it is worth recording, because the two limits
below depend on it. Direct evaluation confirms it: at none of the 400 grid points does
$-\Lmat_{FF}^{-1}$ contain a negative entry.

Differentiating
Eq.~\eqref{eq:R_of_b} logarithmically then gives
\begin{equation}
  \boxed{\ \frac{\mathrm{d}\ln\Robs}{\mathrm{d}\ln\bdep}
  = \ffrac{3} - \ffrac{4}\ }
  \label{eq:log_sensitivity}
\end{equation}

\begin{figure}[tbp]
  \centering
  \includegraphics[width=0.66\textwidth]{fig3_3_ground_fed}
  \caption[Ground-fed fractions and the sensitivity of the observable]{%
  The ground-fed fractions of the two shells and their difference, at the
  benchmark point; the dotted vertical line marks $\bdep^{\mathrm{CRE}}$ there.
  Both fractions rise from zero to unity as the ground-state
  reservoir grows, but they do so at different rates, and their difference,
  which by Eq.~\eqref{eq:log_sensitivity} is exactly
  $\mathrm{d}\ln\Robs/\mathrm{d}\ln\bdep$, vanishes in both limits and
  peaks where the two supplies compete. Note that $\ffrac{3}-\ffrac{4}$ is a
  difference of fractions, not itself a fraction.}
  \label{fig:ground_fed}
\end{figure}

\textbf{The Balmer ratio's sensitivity to the ground-state reservoir is exactly
the difference in ground-fed fraction between the two shells.} This is the
mathematical core of the thesis, and Eq.~\eqref{eq:log_sensitivity} is where the
physics becomes visible.

It vanishes in two limits, and only there. When $\bdep\Zsb{1}\nel \gg
c_m/a_m$ for both shells, both are fed almost entirely from the ground state,
$\ffrac{3} = \ffrac{4} = 1$, and the ratio is blind to the reservoir. When
$\bdep\Zsb{1}\nel \ll c_m/a_m$, both are fed almost entirely by recombination,
$\ffrac{3} = \ffrac{4} = 0$, and the ratio is blind again. \emph{In between, where the two supplies compete, the shells
are fed in different proportion and the sensitivity is largest.}

\subsubsection*{How large, and where}

Equation~\eqref{eq:log_sensitivity} says the sensitivity peaks somewhere between
the two limits. It does not yet say how high, or at what reservoir. Both follow
in closed form, and the answer is more restrictive than one might expect.

Work in the logarithmic reservoir variable $x = \ln(\bdep\Zsb{1}\nel)$. That
combination is the neutral-to-ion ratio itself: $\bdep\Zsb{1}\nel = \ngs/\nion$,
so $x = \ln u$ in the notation Chapter~\ref{ch:results} uses throughout, and the
Saha coefficient cancels out of it entirely. The two chapters work in the same
variable under two names, and it is worth seeing that before the algebra starts.
In that variable Eq.~\eqref{eq:ground_fed_fraction} is exactly a logistic,
\begin{equation}
  \ffrac{m}(x) = \frac{1}{1 + e^{-(x - x_m)}},
  \qquad x_m \equiv \ln\!\left(\frac{c_m}{a_m}\right),
  \label{eq:f_logistic}
\end{equation}

\noindent so the two curves in Figure~\ref{fig:ground_fed} are not merely
similar in shape: they are the \emph{same} curve, of the same width, displaced
along the axis. Every property of the 42-state network enters only through where
each sits. Their separation is
\begin{equation}
  \Delta \equiv x_4 - x_3
  = \ln\!\left(\frac{a_3/a_4}{c_3/c_4}\right)
  = \ln\!\left(\frac{\Robs^{\rm grd}}{\Robs^{\rm rec}}\right),
  \label{eq:logistic_separation}
\end{equation}

\noindent the logarithm of the span between the two asymptotes of
Figure~\ref{fig:R_of_b1}. Equation~\eqref{eq:f_logistic} also makes explicit a
distinction that is easy to lose. \textbf{The switching point of shell $m$ is
$u_m = c_m/a_m$, a ratio of two network coefficients and nothing else.} In the
$\bdep$ coordinate the same point reads
$\bdep = c_m/(a_m\Zsb{1}\nel)$, which carries $\Zsb{1}\nel$ and is therefore a
Saha-normalised label for a quantity that does not depend on the Saha balance.
Equation~\eqref{eq:f_logistic} already says this: it puts the switching point at
$x_m = \ln(c_m/a_m) = \ln u_m$. A plasma may lie far above the Saha--Boltzmann ground-state
population, $\bdep \gg 1$, and still have both shells fed predominantly from
above. The
benchmark point is exactly such a case: $\bdep^{\mathrm{CRE}} = 956$, while the
two shells do not become half ground-fed until $\bdep = 2.73\times10^{3}$ and
$1.91\times10^{4}$ respectively, so $\ffrac{3} = 0.260$ and $\ffrac{4} = 0.048$:
three-quarters of $n=3$ and 95\,\% of $n=4$ arrive by recombination.
A large departure coefficient and a predominantly ground-fed manifold are
different statements, and only the second controls
Eq.~\eqref{eq:log_sensitivity}. Neither should be conflated with calling the
plasma \emph{ionising}, which is properly a statement about the direction of the
net population flux and is not fixed by the magnitude of $\bdep$ alone. Writing $s = x - \tfrac12(x_3+x_4)$ and using
$\ffrac{}(-y) = 1 - \ffrac{}(y)$, the difference $\ffrac{3}-\ffrac{4}$ is
\emph{even} in $s$. Its maximum therefore sits at the midpoint, and takes the
value
\begin{equation}
  \boxed{\ \max\big|\ffrac{3}-\ffrac{4}\big|
  = \tanh\!\left(\frac{|\Delta|}{4}\right)
  = \tanh\!\left[\frac{1}{4}\left|\ln\!\left(\frac{\Robs^{\rm grd}}{\Robs^{\rm rec}}\right)\right|\right]\ }
  \label{eq:peak_sensitivity}
\end{equation}

\noindent where the modulus matters only for the generalisation below: for the
$n_3/n_4$ ratio considered here $\Robs^{\rm grd} > \Robs^{\rm rec}$ and
$\Delta > 0$, but for a pair of shells drawn in the opposite order $\Delta$
changes sign while the left-hand side cannot. The maximum is attained at
$\bdep\Zsb{1}\nel = \sqrt{(c_3/a_3)(c_4/a_4)}$, the
geometric mean of the two switching points. The area is fixed by the same
quantity: integrating Eq.~\eqref{eq:log_sensitivity} over the \emph{whole}
reservoir axis gives
$\int_{-\infty}^{+\infty}(\ffrac{3}-\ffrac{4})\,\mathrm{d}\ln\bdep =
\ln\Robs^{\rm grd} - \ln\Robs^{\rm rec} = \Delta$. The limits are not
decoration: the identity is a sum rule over the full excursion from
purely recombination-fed to purely ground-fed, and an integral over any finite
interval is smaller \emph{in magnitude}. The magnitude matters because the
integrand keeps a fixed sign only while $\Delta > 0$, which holds for the
$(3,4)$ pair; for a pair taken in the opposite order the integrand is negative
throughout and every finite-interval integral is larger, not smaller. Chapter~\ref{ch:results} uses a finite-interval version of
the same integral for a different purpose, to define a mean sensitivity over the
excursion a real step produces, and the two must not be confused. That second relation is only the fundamental
theorem of calculus, and it is quoted here for what it is, a consistency check
tying the two figures together, not an independent result. The figure script
evaluates both and halts if either disagrees with the sampled curve.

Three things follow, and the third is the useful one.

First, the peak sensitivity is set \emph{entirely} by how far apart the two
supply limits lie. Two plasmas with the same span between $\Robs^{\rm rec}$ and
$\Robs^{\rm grd}$ have the same worst-case sensitivity, whatever the details of
their cascade networks.

Second, $\tanh$ is bounded by unity, so
$|\mathrm{d}\ln\Robs/\mathrm{d}\ln\bdep| < 1$ everywhere and for every pair of
shells. \textbf{A shell ratio can never respond to the ground-state reservoir
faster, in relative terms, than the reservoir itself moves.} Integrating a
logarithmic sensitivity whose magnitude is everywhere below unity gives the
finite-excursion statement $|\Delta\ln\Robs| \le |\Delta\ln\bdep|$, so a stale
reservoir cannot be amplified in logarithmic space, only attenuated. The next
subsection shows the same statement holds for the line ratio a spectrometer
actually measures, which is what makes it worth having.

\subsection{The same result for what a spectrometer measures}
\label{sec:emissivity_generalisation}

Everything above is derived for the shell ratio $\Robs = n_3/n_4$, and a
spectrometer does not measure shell populations. It measures line intensities,
which are sums over the dipole-allowed sublevel transitions weighted by their
Einstein coefficients. The gap is worth closing rather than measuring, because
the derivation never used the fact that the weights are unity. It used only
that they are non-negative.

Let $\bm{w}_\alpha$ and $\bm{w}_\beta$ be the weight vectors of the two lines,
carrying $h\nu A$ on the levels each line drains and zero elsewhere. The local
emissivity is $j_\alpha = \bm{w}_\alpha^{\mathsf T}\nF$, and substituting the
two-channel form $\nF = \bm{a}\ngs + \bm{c}\nion$ of
Eq.~\eqref{eq:two_supplies} gives
\begin{equation}
  j_\alpha = A_\alpha\,\ngs + C_\alpha\,\nion ,
  \qquad
  A_\alpha \equiv \bm{w}_\alpha^{\mathsf T}\bm{a},
  \qquad
  C_\alpha \equiv \bm{w}_\alpha^{\mathsf T}\bm{c} ,
  \label{eq:emissivity_two_channel}
\end{equation}

\noindent where $\bm{w}_\alpha$ carries $h\nu A$ on the levels the line drains,
in energy per unit time, and $A_\alpha, C_\alpha \ge 0$ because $\bm{a}$ and
$\bm{c}$ are non-negative elementwise and $\bm{w}_\alpha$ is non-negative. The
logistic form needs the \emph{strict} inequality, since $A_\alpha = 0$ would put
the switching point at infinity, and strictness follows from
$-\Lmat_{FF}^{-1}$ being strictly positive here rather than merely
non-negative; it is confirmed numerically below.
\textbf{That is the same affine form with the same positivity},
so the line ratio
$j_\alpha/j_\beta = (A_\alpha u + C_\alpha)/(A_\beta u + C_\beta)$ is
affine-over-affine in the reservoir exactly as the shell ratio is, and every
result of this section carries over unchanged: the ground-fed fraction of a
\emph{line} is the same unit-width logistic, its switching point is
$C_\alpha/A_\alpha$, the logarithmic sensitivity is the difference of two such
fractions, and the peak is $\tanh(|\Delta_j|/4)$ with
$\Delta_j = \ln[(A_\alpha/A_\beta)/(C_\alpha/C_\beta)]$. Nothing is re-derived;
the substitution $a_m \to \bm{w}^{\mathsf T}\bm{a}$ is the whole of it. The
result holds for any non-negative linear functional of the fast-state vector,
of which a shell sum and an $A$-weighted line are two instances.

\paragraph{Checked, not asserted.} Evaluating
Eq.~\eqref{eq:emissivity_two_channel} at all $400$ grid points with the
pipeline's own resolved $A$-values, $H_\alpha$ from $3s\to2p$, $3p\to2s$ and
$3d\to2p$ and $H_\beta$ from $4s\to2p$, $4p\to2s$ and $4d\to2p$
(\texttt{verify\_emissivity\_generalisation.py}): all four coefficients are
strictly positive everywhere, and maximising $|F_\alpha - F_\beta|$ over $u$
numerically reproduces $\tanh(|\Delta_j|/4)$ to $9\times10^{-16}$, so the bound
is found rather than imposed.

One check is worth more than the others. A constant factor on either weight
vector cancels out of $\Delta_j$, so a photon-counting detector and an
energy-weighted one must return the \emph{same} $\Delta_j$ even though their
emissivities differ by $h\nu$. Measured, the two agree to
$7\times10^{-16}$. \textbf{The displacement that sets the bound is a property of
the atomic network and not of the detector.} The cancellation is exact because
every channel of a given line carries the same $h\nu$ in this model, the levels
being degenerate in $\ell$; with fine structure resolved the two weightings
would differ within a line and the agreement would be approximate rather than
exact.

\paragraph{How far the shell ratio was from the observable.} Over the grid
$\Delta_j$ runs from $0.673$ to $2.196$ against $0.711$ to $2.199$ for the shell
$\Delta$, a median disagreement of \SI{0.134}{\percent} and a worst case of
\SI{5.4}{\percent}; the caps differ by \SI{0.116}{\percent} at the median and
\SI{5.3}{\percent} at worst. At the benchmark, $\Delta_j = 1.9444$ against
$1.9456$, and the cap is $0.4511$ against $0.4514$. The worst case sits at
$\Te = \SI{10}{\electronvolt}$, $\nel = \SI{e12}{\per\cubic\centi\metre}$, the
low-density corner, for the reason Section~\ref{sec:balmer_equivalence} gives:
$4f$ contributes exactly zero to $H_\beta$ while carrying a large share of the
$n=4$ shell, $14/32 = \SI{43.75}{\percent}$ by statistical weight and
\SI{43.66}{\percent} as actually populated at that corner. The two differ
precisely because $\ell$-mixing is weakest there, which is also why the shell
and the line disagree most about what $n=4$ contains.

The shell ratio was therefore a good proxy, and this subsection means the thesis
no longer needs it to be one. \textbf{Section~\ref{sec:balmer_equivalence}'s
measurement of the substitution cost becomes a consistency check on two routes
to the same answer rather than a caveat carried by the analytic results.}

It is worth being precise about what this does and does not bound. It bounds the
error in the \emph{observable}. A diagnostician does not suffer an error in
$\Robs$ directly: they measure $\Robs$ and invert it for $(\Te,\nel)$. Take the
simplest case first, in which $\nel$ has been fixed by an independent
measurement so that only $\Te$ is being inferred. Then the error in the inferred
temperature is
\begin{equation}
  \delta\ln\Te \;=\;
  \frac{\big(\partial\ln\Robs/\partial\ln\bdep\big)_{\Te,\nel}}
       {\big(\mathrm{d}\ln\Rcre/\mathrm{d}\ln\Te\big)_{\nel}}\,\delta\ln\bdep ,
  \qquad \nel \text{ known} .
  \label{eq:inversion_jacobian}
\end{equation}

\noindent Both derivatives are taken at fixed $\nel$, and the relation is a
local linearisation: it is exact only in the limit of small $\delta\ln\bdep$,
and Chapter~\ref{ch:results} uses the exact finite-step form in its place
wherever a number is quoted. It also requires the denominator to be non-zero,
which is not a formality: Section~\ref{sec:inversion_jacobian} finds
$\mathrm{d}\ln\Rcre/\mathrm{d}\ln\Te$ passing through zero inside the operating
range on six of the eight density columns, and there the quotient diverges and
the inversion has no temperature resolution at all.

\textbf{If $\nel$ is not known independently, one ratio is not enough.} Two
unknowns need two observables, and the honest local statement is a Jacobian
rather than a quotient: with two measured ratios $\Robs^{(1)}$ and
$\Robs^{(2)}$,
\begin{equation}
  \begin{pmatrix}\delta\ln\Robs^{(1)}\\[2pt]\delta\ln\Robs^{(2)}\end{pmatrix}
  = \bm{J}\begin{pmatrix}\delta\ln\Te\\[2pt]\delta\ln\nel\end{pmatrix},
  \qquad
  \begin{pmatrix}\delta\ln\Te\\[2pt]\delta\ln\nel\end{pmatrix}
  = \bm{J}^{-1}\begin{pmatrix}\delta\ln\Robs^{(1)}\\[2pt]\delta\ln\Robs^{(2)}\end{pmatrix},
  \label{eq:inversion_2x2}
\end{equation}

\noindent where $\bm{J}$ is the Jacobian of $(\ln\Robs^{(1)},\ln\Robs^{(2)})$
with respect to $(\ln\Te,\ln\nel)$, and the inversion exists only where
$\det\bm{J} \ne 0$. That condition is the classical failure mode of
two-ratio diagnostics rather than a technicality, because the two families of
contours in the $(\Te,\nel)$ plane can run nearly parallel over wide regions.
To parallel Eq.~\eqref{eq:inversion_jacobian}, a reservoir displacement
$\delta\ln\bdep$ enters as
$(\delta\ln\Te, \delta\ln\nel)^{\mathsf T} = \bm{J}^{-1}
(\partial\ln\Robs^{(1)}/\partial\ln\bdep,\,
\partial\ln\Robs^{(2)}/\partial\ln\bdep)^{\mathsf T}\,\delta\ln\bdep$,
and the amplification is governed by the conditioning of $\bm{J}$
rather than by a single quotient. Chapter~\ref{ch:intro} already noted that one
ratio is one constraint on two unknowns and that practice supplies a second;
Eq.~\eqref{eq:inversion_jacobian} is what that second constraint buys.
This thesis measures the numerator and does not measure $\bm{J}$, so no claim
about a two-parameter inversion is made anywhere in it.

The bound constrains the numerator only. Where the ratio is itself
weakly sensitive to temperature, which is where Balmer diagnostics are known
to be least useful, a bounded error in $\Robs$ can still become a large error
in the inferred $\Te$. The inversion Jacobian in the denominator is supplied in
Chapter~\ref{ch:results}; without it, no statement about diagnostic error
follows from this section.

These closed forms also locate the benchmark point on the curve rather than
leaving it as an anchor of unstated significance. Its sensitivity is
$\ffrac{3}-\ffrac{4} = 0.212$, some $47\,\%$ of the maximum $0.4514$, which is
attained nearly one decade higher at $\bdep = 7.2\times10^{3}$. The benchmark is
therefore on the rising flank and not at the worst case, which is the honest
reading of every error reported against it in
Chapter~\ref{ch:results}. Two identities follow from
Eq.~\eqref{eq:ground_fed_fraction} and fix the marker in
Figure~\ref{fig:R_of_b1} in terms of the fractions of
Figure~\ref{fig:ground_fed},
\begin{equation}
  \Robs = \Robs^{\rm grd}\,\frac{\ffrac{4}}{\ffrac{3}}
        = \Robs^{\rm rec}\,\frac{1-\ffrac{4}}{1-\ffrac{3}},
  \label{eq:R_from_fractions}
\end{equation}

\noindent both giving $\Robs = 0.7778$ at the benchmark. They are algebraic
rearrangements rather than new physics, but they are the cheapest available
cross-check between the two figures and the figure script enforces them.

Third, the attenuation is severe unless the span is large. For $\Delta \ll 1$
the peak is $\Delta/4$: a diagnostic whose two supply limits differ by only
$10\,\%$ has a worst-case log-sensitivity of $0.024$. Sensitivity to the hidden
variable is bought with span, and $\Robs = n_3/n_4$ has span because $n=3$ and
$n=4$ draw on the two supplies in genuinely different proportion. That is why
this ratio is diagnostically useful and also why it is exposed.

Note that this is the logarithmic sensitivity, not the absolute derivative
$\partial\Robs/\partial\bdep$ of Eq.~\eqref{eq:dR_db}. The latter does
\emph{not} vanish as $\bdep \to 0$; it tends to
$\Zsb{1}\nel(a_3c_4-a_4c_3)/c_4^2$, which is finite. The two limits are
statements about the \emph{fractional} response of the observable, which is what
a diagnostic error is.

The two supply channels are Fujimoto's, who names the first term of
Eq.~\eqref{eq:two_supplies} the \emph{recombining plasma component} and the
second the \emph{ionising plasma component}~\cite[Sec.~4.1]{Fujimoto2004}.
Naming the components that way is his and is standard. \textbf{Identifying which
component dominates one shell with the state of the plasma is a different claim,
and this thesis does not make it.} Call a shell \emph{ground-fed} when
excitation out of the ground state supplies most of its population and
\emph{recombination-fed} when capture from the continuum does. That is a
statement about one level's supply budget. Whether the plasma is \emph{ionising}
is a statement about the sign of the net stage-to-stage flux, as
Section~\ref{sec:ground_fed_fraction} already insisted, and a level can be
predominantly ground-fed in a plasma that is recombining on balance. The
crossover Eq.~\eqref{eq:log_sensitivity} singles out is the ground-fed one:
between the two limits both supply channels contribute, the two shells draw on
them in different proportion, and the sensitivity is largest. That crossover
is a statement about the switching points $c_m/a_m$ of the network, and where it
falls in $(\Te,\nel)$ is measured rather than assumed; identifying it with any
named divertor operating regime is a separate claim, requiring densities this
section does not supply, and is not made here.
Chapter~\ref{ch:results} maps $\ffrac{3} - \ffrac{4}$ over the operating range
and shows the diagnostic error tracking it.

It is worth recording what Eq.~\eqref{eq:R_of_b} is not. It is not a fit, an
expansion, or an approximation: it is an exact rewriting of
Eq.~\eqref{eq:qss_solution}, valid wherever $\Lmat_{FF}$ is invertible.
Nor is it new. It is the standard population-coefficient decomposition of
collisional--radiative theory, introduced by Bates, Kingston and
McWhirter~\cite{Bates1962} % [KEY REQUIRED: Proc. Roy. Soc. A 267, 297 (1962)]
and written by Fujimoto and
McWhirter~\cite{FujimotoMcWhirter1990} % [KEY REQUIRED: Phys. Rev. A 42, 6588 (1990)]
in the Saha-normalised form
$b(p) = r_0(p) + r_1(p)\,b(1)$, with $r_0$ the recombining and $r_1$ the
ionising population coefficient. Substituting
$b(p) = n_p/[\Zsb{p}\nel\nion]$ into that expression and clearing the
normalisation gives the correspondence
\begin{equation}
  a_p = r_1(p)\,\frac{\Zsb{p}}{\Zsb{1}},
  \qquad
  c_p = r_0(p)\,\Zsb{p}\,\nel ,
  \label{eq:fujimoto_correspondence}
\end{equation}

\noindent both sides dimensionless. The two forms are the same linear solution
in different coordinates; what is done here is to ask what the ratio of two such
populations depends on, which Eq.~\eqref{eq:log_sensitivity} answers.

Section~\ref{sec:fujimoto} tests the decomposition against Fujimoto's published
population coefficients. A second test of the same decomposition is attempted
against the ADAS effective coefficients~\cite{Summers2006}, and it fails.
Because $r_0$ and $r_1$ are published as curves rather than tables, that test
cannot be made level by level: $\bm{a}$ and $\bm{c}$ are instead summed against
the ionisation rates to form the effective stage-to-stage coefficients, which
\emph{are} tabulated, and compared with SCD96 and ACD96. The comparison,
Gate~D of Chapter~\ref{ch:validation}, disagrees at every grid point, by a
systematic factor between $9$ and $65$ that peaks near
$\Te \approx \SI{2}{\electronvolt}$, and only its ionisation half is
implemented, the recombination table being read and never used.
Chapter~\ref{ch:validation} reports that failure and its diagnosis rather than
repairing it. Two readings are open. ADAS tabulates SCD as the \emph{ionising}
coefficient, whereas the sum formed here runs over the full steady state, so
ionisation out of recombination-fed Rydberg levels enters the model's
coefficient and belongs under ACD; and this model holds the ion as a fixed
reservoir with no forty-fourth state (Section~\ref{sec:two_references}), so the
two quantities are not defined on the same system. Neither reading has been
demonstrated, and no agreement with ADAS is claimed anywhere in this thesis.

The contraction is severe in any case, since forty-two coefficients are
collapsed into one and many incorrect $\bm{a}$ would contract correctly, but it
is not merely a rate check either:
$a_p$ and $c_p$ enter the two coefficients with different weights, so an
error in the \emph{split} between the two supplies displaces the ionisation and
recombination coefficients in opposite directions, which no single-channel
comparison detects. Chapter~\ref{ch:results} verifies the superposition
$\nF = \bm{a}\ngs + \bm{c}\nion$ to a relative error of
$3\times10^{-14}$.

\subsection{Two reference states, and they are not the same}

Equation~\eqref{eq:R_of_b} shows that the QSS ratio is a function of three
arguments, not two. This gives two distinct things a measurement can be compared
against.

\paragraph{The QSS reference.} $\Rqss(\Te, \nel;\,\bdep(t))$, evaluated at the
\emph{instantaneous} departure coefficient. This is what adiabatic elimination
claims: the excited manifold is in steady state with whatever reservoir is
present at that moment. Its accuracy is governed by
Section~\ref{sec:qss}'s validity condition.

\paragraph{The CR-equilibrium reference.}
\begin{equation}
  \Rcre(\Te, \nel)
  \equiv \Rqss\!\left(\Te, \nel;\,\bdep^{\mathrm{CRE}}(\Te,\nel)\right),
  \label{eq:cre_ratio}
\end{equation}

\noindent where $\bdep^{\mathrm{CRE}}$ is obtained by additionally setting
$\mathrm{d}\nS/\mathrm{d}t = 0$ in Eq.~\eqref{eq:reduced} and solving for the
steady ground-state population. This \emph{is} a function of $(\Te,\nel)$ alone
and can therefore be tabulated.

The distinction is easy to miss because both are called ``the QSS solution'' in
common usage, and in a steady plasma they are the same number. They separate
whenever the ground-state population has not finished settling, which, in a
divertor disturbed by an edge-localised mode every few milliseconds, is most of
the time.

Two remarks on terminology, both of which matter for what follows. First,
$\bdep = 1$ denotes complete thermodynamic equilibrium, whereas
$\bdep^{\mathrm{CRE}}$ denotes the steady state of the collisional--radiative
system, which is a different and much weaker condition; where excitation out of
the ground state dominates and the plasma is optically thin,
$\bdep^{\mathrm{CRE}}$ is very much larger than unity. Note that a large
$\bdep^{\mathrm{CRE}}$ does not by itself make the shells ground-fed: the
benchmark has $\bdep^{\mathrm{CRE}} = 956$ and both shells fed predominantly
from above, which is the distinction
Section~\ref{sec:ground_fed_fraction} draws. Second,
this model holds the ion density fixed (Section~\ref{sec:rate_equation}), so
the quantity relaxing on $\tauslow$ is the ground-state population against a
prescribed reservoir, not a self-consistent ionisation balance. What that
restriction costs is assessed in Chapter~\ref{ch:limits}.

\subsection{Two error measures}

Each reference defines an error:
\begin{align}
  \epsQSS(t) &= \left|\frac{\Robs(t)}
                     {\Rqss\!\left(\Te,\nel;\,\bdep(t)\right)} - 1\right|,
  \label{eq:eps_qss}\\[4pt]
  \epsCRE(t) &= \left|\frac{\Robs(t)}
                     {\Rcre\!\left(\Te,\nel\right)} - 1\right|.
  \label{eq:eps_cre}
\end{align}

$\epsQSS$ asks whether the excited states have kept up with the reservoir they
actually see. $\epsCRE$ asks whether the tabulated answer is right. The first is
the question the literature worries about; the second is the one a
diagnostician suffers.

\textbf{They are not the same question, and Chapter~\ref{ch:results} shows they
have very different answers.} Naming them apart is the single most important
definitional step in this thesis: conflating them makes it possible to measure
the second and report it as the first.

\subsection{What a two-parameter table must assume}

It remains to say precisely what is being assumed when a measured ratio is
inverted against a table indexed on $(\Te,\nel)$.

Equation~\eqref{eq:R_of_b} establishes that $\Robs = \Robs(\Te,\nel,\bdep)$.
A two-parameter table must therefore do one of three things: adopt a closure
$\bdep = \bdep(\Te,\nel)$; operate in a regime where one supply dominates both
shells, so that Eq.~\eqref{eq:dR_db} is negligible and the third argument does
not matter; or hold some other quantity fixed that determines $\bdep$
implicitly. Equilibrium ionisation balance is a closure, not a mathematical
necessity, and it is important not to overstate the point.

It is important to be precise about which practice is being criticised, because
the modern effective-coefficient databases are not themselves at fault.
ADAS~\cite{Summers2006}, for instance, tabulates photon emissivity coefficients
separately for the excitation-driven and recombination-driven components. Used
as intended, those two components are weighted by $\ngs$ and $\nion$
respectively, which is exactly the information $\bdep$ carries; the database
does not impose a closure at all. Indeed the excitation/recombination split of
Eq.~\eqref{eq:two_supplies} is conceptually the same decomposition.

The assumption enters at the \emph{inversion} step, when a measured ratio is
mapped to $(\Te,\nel)$ through a two-parameter lookup. Such a lookup returns a
single ratio per $(\Te,\nel)$ pair and therefore requires the two components to
have been combined in a fixed proportion; that is, it requires a closure for
$\bdep$. The usual choice is $\bdep = \bdep^{\mathrm{CRE}}$, obtained by solving
the collisional--radiative system to complete steady state. A ratio inverted this
way carries the assumption that the plasma's ground-state population has settled
at its steady value for the local conditions: not an assumption about how quickly
excited states equilibrate, but an assumption that the plasma has \emph{finished
changing}.

The second possibility is not merely hypothetical, and it is the reason
Eq.~\eqref{eq:dR_db} is worth having. Where one supply dominates both shells,
the table is safe regardless of what the ground state is doing.
Chapter~\ref{ch:results} maps where in the divertor operating range that is the
case, and where it is not.

The remainder of this thesis quantifies what the assumption costs.

%-----------------------------------------------------------------------------
%-----------------------------------------------------------------------------
\section{A caveat: the modes are not orthogonal}
\label{sec:transient_growth}

Section~\ref{sec:two_clocks} presented the eigen-decomposition as a change of
coordinates in which the 43 coupled equations become 43 independent decays. That
is exact. But a change of coordinates is only as convenient as the coordinates
are well behaved, and these are not.

The reason is a physical asymmetry running through the whole matrix. Radiative
decay carries atoms downward at up to $\SI{6.3e8}{\per\second}$; electron-impact
excitation carries them upward at $C\nel \sim \SI{e4}{\per\second}$ at the
benchmark density. Downward and upward coupling therefore differ by nearly five
orders of magnitude, and $\Lmat$ is very far from symmetric. That asymmetry does
not by itself establish non-normality, for the reason given in
Section~\ref{sec:two_clocks}: a non-symmetric matrix may still be normal. The
non-normality is established separately, by measurement, and the numerical
abscissa is positive at all $400$ grid points.

The consequence is worth stating carefully, because it is easy to overstate.
Every eigenvalue has negative real part, so every mode decays and any
perturbation eventually vanishes; that much the spectrum does guarantee.
What it does not guarantee is that the perturbation decays \emph{monotonically}.
When eigenvectors are nearly parallel, a state of modest size can be the sum of
two much larger, nearly cancelling modal contributions. If those two modes decay
at different rates the cancellation is progressively spoiled, and the state can
grow, sometimes by orders of magnitude, before the decay of the individual
modes finally wins. The standard diagnostic for whether this is possible at all
is the numerical abscissa
\begin{equation}
  \mu(\Lmat) = \lambda_{\max}\!\left(\tfrac{1}{2}\left[\Lmat + \Lmat^{\mathsf T}\right]\right),
  \label{eq:numerical_abscissa}
\end{equation}

\noindent the largest eigenvalue of the symmetric part of $\Lmat$, which gives
the initial slope of $\|\nvec(t)\|$ for the worst-case perturbation. For this
matrix $\mu(\Lmat) > 0$ at every grid point, so transient growth is not merely
possible but guaranteed for \emph{some} disturbance.

Whether it occurs for the disturbance and the observable that matter here is a
different question, and one the spectrum cannot answer. $\mu(\Lmat) > 0$ is a
statement about the worst case over all perturbation directions in the Euclidean
norm, and other norms carry their own logarithmic norms; the Balmer
ratio is one particular scalar function of the state, and a temperature step is
one particular disturbance. A quantity may be entirely well behaved even when
the operator that governs it admits growth in principle. This is why
Section~\ref{sec:qss} was careful to label Eq.~\eqref{eq:qss_validity} an
estimate rather than a bound: the eigenvalue argument describes the system after
the transients have settled, and the transients themselves are outside its
reach.

Chapter~\ref{ch:results} therefore settles the question the only way it can be
settled, by integrating the full 43-state system through a temperature step
and measuring what the Balmer ratio does, rather than by inferring it from the
spectrum.

%-----------------------------------------------------------------------------
\section{Summary}
\label{sec:theory_summary}

Starting from a population balance on each of 43 levels, this chapter assembled
the collisional--radiative rate equation $\dot{\nvec} = \Lmat\nvec + \Svec\nion$
and checked the assembly against particle conservation, which holds to
$3.6\times10^{-11}$ across the operating range.

Diagonalising $\Lmat$ revealed a spectrum with one gap and then a continuum: a
single slow mode at $\tauslow = \SI{22.7}{\micro\second}$, the relaxation of the
ground-state population, sitting four orders of magnitude above a dense band of
fast modes whose slowest member is $\taurel = \SI{2.28}{\nano\second}$, the
rearrangement of the excited manifold. The separation
$\Msep = \tauslow/\taurel$ is 9982 at the benchmark point and does not fall below
$86.8$ at any of the 400 grid points, though, as
Section~\ref{sec:two_clocks} records, that floor sits at a point the transient
analysis excludes, and the four-orders figure quoted above is the benchmark
value rather than a property of the grid, over which $\Msep$ spans two to nine
orders.

That separation licenses the quasi-steady-state approximation: eliminate the
excited populations algebraically and solve only for the slow subsystem. The
elimination is the Bodenstein steady-state approximation applied to 42 coupled
intermediates at once, and it produces the Schur complement $\Omqss$, the
effective ionisation rate that a plasma modeller using QSS constructs and the
algebraic form of the effective coefficients that spectroscopic inversions rest
on. Its validity condition, $\Msep \gg 1$, is satisfied everywhere on the grid
and comfortably, subject to the two qualifications of
Section~\ref{sec:qss}: $\Msep$ compares rates where
Chapter~\ref{ch:results} measures a distance, and the spectrum of a non-normal
operator bounds asymptotic decay rather than finite-time error.

The chapter's central result is that this is not the end of the argument. The
QSS ratio
\begin{equation*}
  \Rqss(\bdep) = \frac{a_3\,\bdep\,\Zsb{1}\,\nel + c_3}
                      {a_4\,\bdep\,\Zsb{1}\,\nel + c_4}
\end{equation*}

\noindent is a function of \emph{three} arguments, not two. The third is the
ground-state departure coefficient $\bdep$, which is dynamical: it carries the
plasma's history and takes a time $\tauslow$ to settle. Its logarithmic
sensitivity is exactly
\begin{equation*}
  \frac{\mathrm{d}\ln\Robs}{\mathrm{d}\ln\bdep} = \ffrac{3} - \ffrac{4},
\end{equation*}

\noindent the difference in ground-fed fraction between the two shells, which
vanishes when either supply dominates both shells and is largest where the two
compete. A table indexed on $(\Te,\nel)$ alone must therefore adopt a closure for
$\bdep$, and the closure adopted in practice is complete collisional--radiative
equilibrium.

Two error measures follow and must not be conflated. $\epsQSS$ asks whether the
excited states have kept up with the reservoir they actually see; $\epsCRE$ asks
whether the tabulated answer is right. The first is the question the literature
worries about. The second is the one a diagnostician suffers.

Everything in this chapter has been an argument about equations. Whether the
equations describe hydrogen is a separate question, and Chapter~\ref{ch:validation}
addresses it: what would a wrong model fail, and does this one fail it?
